% registry: RFB
\methodhead{RFB}{Robert, Fogue and Bahuaud (1994): integral criterion}\label{sheet:RFB}

\block{Source} J.-L.~Robert, M.~Fogue, J.~Bahuaud, in \emph{Automation in Fatigue and Fracture:
Testing and Analysis}, ASTM STP 1231 (1994) 369--387 \cite{robert1994}, Eqs.~(8), (15)--(25), (27),
Appendix; Tables~1--3.

\block{Idea and scope} Integral criterion. On each plane a damage indicator combines linearly the
largest projected shear stress amplitude, the normal stress amplitude and the normal mean stress;
its quadratic mean over all planes through the point is compared with the fully reversed axial
strength. Three weights calibrate the integral to reversed tension, reversed torsion and pulsating
tension. Proposed for periodic multiaxial stress states, and for random ones after counting, at a
given life; valid for $0.577<t_{-1}/s_{-1}<0.866$.

\block{Inputs}
\begin{center}\small
\begin{tabular}{@{}llp{100mm}@{}}
\toprule
Symbol & Unit & Meaning\\\midrule
$\mathbf h$, $S$ & -- & plane normal and the unit sphere ($\lvert S\rvert=4\pi$)\\
$\psi$ & rad & direction of a line in the plane $\mathbf h$\\
$\tau_{h\psi,a}$ & MPa & amplitude of the shear stress projected on line $\psi$\\
$\sigma_{hh,a}$, $\sigma_{hh,m}$ & MPa & normal stress amplitude and mean on the plane\\
$E_h$ & MPa & damage indicator of the plane, Eq.~\eqref{eq:rfb-eh}\\
$E$ & -- & validity index, Eq.~\eqref{eq:rfb-main}; $I$ in Eq.~(27)\\
$s_{-1}=\fm$, $t_{-1}=\tm$ & MPa & fully reversed axial and torsion strengths (fatigue limits in the repository)\\
$s_0$ & MPa & pulsating ($R=0$) axial strength, maximum stress\\
$\sigma_u$ & MPa & ultimate tensile strength (estimate of $s_0$)\\
$k=s_{-1}^2/t_{-1}^2$, $D$ & --, -- & limit-ratio square and discriminant of Eq.~\eqref{eq:rfb-ab}\\
$a$, $b$, $d$ & -- & weights of the shear amplitude, normal amplitude and normal mean\\
$c$ & -- & further weight of the paper's general indicator, set to 0 by Eq.~(25)\\
$\sigma_A$, $\sigma_M$; $\tau_A$, $\tau_M$ & MPa & axial and torsional amplitude and mean on the Haigh curves\\
$\sigma_a$, $\tau_a$ & MPa & bending and torsion amplitudes of the worked example\\
\bottomrule
\end{tabular}
\end{center}

\block{Formulation} With $c=0$ (Eq.~25), the plane indicator and the failure condition
(Eqs.~8, 15--16) are
\begin{gather}
E_h=a\max_{\psi}\tau_{h\psi,a}+b\,\sigma_{hh,a}+d\,\sigma_{hh,m},
\label{eq:rfb-eh}\\
E=\frac{1}{s_{-1}}\sqrt{\frac{1}{4\pi}\int_S E_h^2\,dS}=1 .
\label{eq:rfb-main}
\end{gather}
Reversed tension, reversed torsion and pulsating tension at the limits ($\sigma_A=\sigma_M=s_0/2$)
give, in this order, the conditions of Eqs.~(20), (21) and (24),
\begin{equation}
2a^2+3b^2+4ab=15,\qquad 6a^2+4b^2+7ab=15\,k,\qquad 15+3d^2+6bd+4ad=60\,\frac{s_{-1}^2}{s_0^2},
\label{eq:rfb-cal}
\end{equation}
solved in the Appendix as
\begin{gather}
D=225\left(1-\frac{8(k-3)^2}{25}\right),\qquad b^2=\frac{15-\sqrt D}{2},\qquad a^2=\frac{12k-21+b^2}{2},
\label{eq:rfb-ab}\\
d=\frac{-(3b+2a)+\sqrt{(3b+2a)^2+3\,(60\,s_{-1}^2/s_0^2-15)}}{3}.
\label{eq:rfb-d}
\end{gather}
Real weights exist only for $4/3<k<3$, i.e.\ $0.577<t_{-1}/s_{-1}<0.866$. The Haigh curves in tension
(Eq.~19) and torsion (Eq.~23) follow as
\begin{equation}
15\,s_{-1}^2E^2=15\,\sigma_A^2+3d^2\sigma_M^2+2d(3b+2a)\,\sigma_A\sigma_M ,\qquad
15\,s_{-1}^2E^2=15\,k\,\tau_A^2+4d^2\tau_M^2 ,
\label{eq:rfb-haigh}
\end{equation}
where the mean shear acts only through the normal stresses of oblique planes.

\block{Constants and calibration} $a,b,d$ in closed form from $s_{-1}$, $t_{-1}$, $s_0$; nothing is
fitted statistically. The paper uses the strengths at the target life $N$; the repository uses the
workbook fatigue limits \texttt{Sigm1}, \texttt{Taum1} and \texttt{Sig0}, or, without
\texttt{Sig0}, the estimate $s_0=4s_{-1}\sigma_u/(s_{-1}+2\sigma_u)$ of B\"ohme, Papuga and Lange
(2026), Eq.~(1), from \texttt{Su}. Only the master-curve constants are fitted, on the uniaxial
$R=-1$ design tests.

\block{Life} The paper applies the criterion with the strengths at life $N$. The repository returns
the validity index $E$ computed with the fatigue limits and reads finite lives on the master curve
of $E$ (Section~\ref{sec:mastercurve}).

\block{Repository choices}
\begin{itemize}
\item Misprints corrected: Eq.~(16) prints $1/\pi$ instead of the $1/(4\pi)$ required by Eqs.~(8) and
  (18)--(24); the Appendix omits the $\tfrac12$ under the root of $a$, required by Eqs.~(20) and (24);
  Eq.~(18) is read with $E^2$ where $E$ is printed.
\item Fatigue limits replace the strengths at life $N$ (see Life).
\item $t_{-1}/s_{-1}$ outside $(0.577,0.866)$ raises \texttt{MaterialScopeError}: the model is skipped
  on that material (eligibility \texttt{out\_of\_scope}), e.g.\ AISI316L with $t_{-1}/s_{-1}=0.517$.
\item For a single harmonic, $\max_\psi\tau_{h\psi,a}$ is the radius of the minimum circumscribed
  circle of the shear path on the plane.
\item Sphere average on a 48$\times$192 Gauss--Legendre grid.
\end{itemize}

\block{Worked example} Steel of Table~2: $s_{-1}=361$, $t_{-1}=228$, $s_0=600\MPa$. In-phase reversed
bending $\sigma_a=250\MPa$ with torsion $\tau_a=120\MPa$: $E=0.8698$.

\block{Verification} Checked for the three materials of Tables~1--3
($s_{-1},t_{-1},s_0$ = 423, 287, 712; 361, 228, 600; 695, 415, 1040\MPa): the Appendix solves
Eqs.~\eqref{eq:rfb-cal}, and the printed $a$ violates Eq.~(20); the three calibration points give
$E=1$; the Haigh curves Eqs.~\eqref{eq:rfb-haigh}, to $4\cdot10^{-4}$; the validity range. Tables~1--3:
the validity parameter $I$ is reproduced to print rounding (0.006) on 52 of 59 tests. Strict
\texttt{xfail}s: Table~1 rows 1, 2, 5, 6, 7, 12 (biaxial tension--pressure, 0.006--0.015 off,
unresolved) and Table~3 row 11 ($\delta=45^\circ$, printed 1.04 against 1.073 from Eq.~(16), which
breaks the paper's own phase series).
